%------------------------------------------------------------------------------------------
\subsection{Group I (with c.c.)}
\subsubsection{Row I}
Only four-$\rho$ terms, all $\propto\log(\vee/\Lambda)$ with $A=P$: \textbf{$B=0$}.

\subsubsection{Row II}\label{sec:I.II}
\paragraph{Part 1 (two-$\rho$).} Your bracket is $\big[\frac{11}{3}+\log\frac{\vee}{k^+}-2\log\frac{\vee}{\Lambda}+\log(\frac{\vee}{k^+}-1)\big]$, and $\log\frac{\vee}{k^+}+\log(\frac{\vee}{k^+}-1)=2\log\frac{\vee}{k^+}+\log(1-\frac{k^+}{\vee})$. The prefactor $\frac{g^4N_c}{32\pi^4k^+}$ is $\frac{\pi N_c}{2}$, so
\begin{align}
A_{\rm I.II,1}&=-\pi N_c\,\NN\int\KK(x'-w')\cdot\KK(x-w)\,B(x,w)\,\big[U^{ab'}(x')-U^{ab'}(w')\big]\rho^{b'}(x')\big[U^{ab}(w)-U^{ab}(x)\big]\rho^b(x)+{\rm c.c.},\notag\\
B_{\rm I.II,1}&=+\pi N_c\,\NP\int\KK(x'-w')\cdot\KK(x-w)\,B(x,w)\,\big[U^{ab'}(x')-U^{ab'}(w')\big]\rho^{b'}(x')\big[U^{ab}(w)-U^{ab}(x)\big]\rho^b(x)+{\rm c.c.},
\end{align}
and $P_{\rm I.II,1}=0$. As in the previous note (fix F6), the bracket must be $B(x,w)=\frac1\pi\int_z[K(x,w,z)-K(w,w,z)]=\frac1{\bar\epsilon}-\log(4\pi^2\mu^2(x-w)^2)$ instead of $-\frac2\epsilon-2\gamma-\log\frac{(x-w)^2\mu^2}{4}$.

\paragraph{Part 2 (three-$\rho$, momentum space).} In your bracket (with $\mathbf k$ the Fourier variable inside the bracket), Table~\ref{tab:L} gives, term by term:
\begin{align*}
&\log\tfrac{\Lambda}{k^+}\,\tfrac{k^2p^i-k\cdot p\,k^i}{(k-p)^2}&&\to\ A:-\tfrac{k^2p^i-k\cdot p\,k^i}{(k-p)^2},\quad B:+\tfrac{k^2p^i-k\cdot p\,k^i}{(k-p)^2},\\
&-2k^i\log\tfrac{\vee}{\Lambda}&&\to\ A:-2k^i,\\
&-2k^i\log\tfrac{k^+p^2}{\vee k^2+k^+p^2}&&\to\ B:+2k^i,\\
&\tfrac{(2p^2-2k\cdot p+k^2)k^i+k^2p^i}{(k-p)^2}\log\big[1+(\tfrac{\vee}{k^+}-1)\tfrac{k^2}{(k-p)^2}\big]&&\to\ B:+\tfrac{(2p^2-2k\cdot p+k^2)k^i+k^2p^i}{(k-p)^2},\\
&-\tfrac{p^2k^i+k^2p^i}{(k-p)^2}\log\tfrac{\vee-k^+}{\Lambda}&&\to\ A:-\tfrac{p^2k^i+k^2p^i}{(k-p)^2},\\
&-k^i\log\tfrac{\vee}{k^+}&&\to\ B:-k^i.
\end{align*}
So the $\log(\vee/k^+)$ bracket is
\begin{equation}
B\text{-bracket}=\frac{2k^2p^i+(3p^2-5k\cdot p+2k^2)k^i}{(k-p)^2},\qquad A\text{-bracket}=-\frac{2k^2p^i+(p^2-k\cdot p)k^i}{(k-p)^2}-2k^i,
\end{equation}
and their sum (the coefficient of $\log\vee$) is $\frac{2(k-p)^2k^i}{(k-p)^2}-2k^i=0$. Hence $P_{\rm I.II,2}=0$ and
\begin{align}
B_{\rm I.II,2}=-A_{\rm I.II,2}={}&+\frac{i}{2}\,\NP\int\KK^i(x'-w')\,\mathfrak F^i(w;x,y)\,\big[U^{ab'}(x')-U^{ab'}(w')\big]\rho^{b'}(x')\notag\\
&\times\Big[U^{ah}(w)f^{hbc}\{\rho^b(x),\rho^c(y)\}-f^{abc}U^{bd}(x)U^{ce}(y)\{\rho^d(x),\rho^e(y)\}\Big]+{\rm c.c.},
\end{align}
with $\mathfrak F^i=\mathfrak F^i_\alpha+\mathfrak F^i_\beta$ the position-space kernel of the three-$\rho$ note (App.~A there).

\textbf{Fix F4 does not change $B$.} The previous note removed the $\alpha$ part $\mathfrak F_\alpha$ from the $\log(\vee/\Lambda)$ terms of the $U(w)\mathbb A^{(3)}$ sub-terms ($A\to A-\alpha$). The region-$P$ comparison (Sec.~\ref{sec:checks}) shows that the same $\alpha$ part must also be removed from $P$ ($P=0\to-\alpha$). So the $\alpha$ part multiplies $\log(\vee/\Lambda)$ only, and $B=P-A$ is unchanged: the $\log(\vee/k^+)$ terms of Row II are correct as written.

\subsubsection{Row III}\label{sec:I.III}
Prefactor $\frac{ig^4f^{abd}}{16\pi^5k^+}\to i\,f^{abd}$ and colour operator
\[
\mathcal O_{\rm III}=\big[U^{ab'}(x')\rho^{b'}(x')-U^{ab'}(w')\rho^{b'}(x')\big]\big[U^{de}(x)U^{bc}(z)\rho^e(x)\rho^c(y)-U^{de}(x)U^{bc}(y)\rho^e(x)\rho^c(y)\big].
\]
With $a=(x-w)^2$, $b=(x-z)^2$, $X=x'-w'$, $Y=y-z$, $Z_w=z-w$, the kernels of your parts are
\begin{align*}
&K_a=\frac{X\cdot Y\,(a+b)}{X^2Y^2(b-a)^2}=-\KK(x'-w')\cdot\KK(y-z)\,\Phi(x),\qquad K_b=\KK(x'-w')\cdot\KK(x-w)\,K(x,y,z),\\
&K_c=\KK(x'-w')\cdot\KK(x-w)\,\KK(y-z)\cdot\KK(z-w),\qquad K_h=\KK(x'-w')\cdot\KK(z-w)\,K(x,y,z),\\
&K_d=\frac{X\cdot Y\,Z_w\cdot(x-z)\,a}{X^2Y^2Z_w^2(a-b)^2},\quad K_e=\frac{X\cdot Y\,Z_w\cdot(x-w)\,b}{X^2Y^2Z_w^2(a-b)^2},\quad K_f=\frac{X\cdot Y}{X^2Y^2Z_w^2},
\end{align*}
and $g_1,\dots,g_4$ the four terms with $\frac1{a-b}$ in your part 2 (multiplying $L_o(a,b)$, in order: $+\frac{X\cdot Z_w\,Y\cdot(x-w)}{(a-b)}$, $-\frac{X\cdot(x-w)\,Y\cdot Z_w\,b}{(a-b)a}$, $-\frac{X\cdot Z_w\,Y\cdot(x-z)\,a}{(a-b)b}$, $+\frac{X\cdot(x-z)\,Y\cdot Z_w}{(a-b)}$, all over $X^2Y^2Z_w^2$). Your longitudinal functions give (units $i f^{abd}$, Table~\ref{tab:L}):
\begin{center}\small
\begin{tabular}{@{}l l c c@{}}
part & your term & $A$ & $B$\\ \hline
1 & $-[\log\frac{\vee}{\Lambda}+\log(1-\frac{k^+}{\vee})]\,K_a$ & $-K_a$ & 0\\
1 & $+L_o(a,b)\,K_a$ & 0 & $+K_a$\\
1 & $-[\log\frac{\vee}{k^+-\Lambda}-\log\frac{\vee}{\Lambda}]\,K_a$ & $+K_a$ & $-K_a$\\
1 & $+\log\frac{\vee}{\Lambda}\,K_b$ & $+K_b$ & 0\\
1 & $-2\,L_{o'}(a,b)\,K_b$ & 0 & $-2K_b$\\
2 & $-2\log\frac{\vee}{\Lambda+k^+}\,K_c$ & 0 & $-2K_c$\\
2 & $+2[\log\frac{\vee}{\Lambda}+\log(1-\frac{k^+}{\vee})]\,(K_d-K_e)$ & $2(K_d-K_e)$ & 0\\
2 & $+2L_o(a,b)\,[g_1+g_2+g_3+g_4+K_e-K_d]$ & 0 & $2(g_1+g_2+g_3+g_4)+2(K_e-K_d)$\\
3a & $[\log\frac{\vee}{\Lambda}+\log\frac{k^+-\Lambda}{\vee}]\,(2K_d-2K_e-2K_f)$ & $2(K_d-K_e-K_f)$ & $-2(K_d-K_e-K_f)$\\
3b & $\mathcal A$, $\log\frac{[k^+(x-w)-(k^+-\Lambda)(x-z)]^2}{[k^+(x-w)-\Lambda(x-z)]^2}$ terms & 0 & 0\\
\end{tabular}
\end{center}
Two identities simplify the sums: $4(K_d-K_e)-2K_f=-2K_a$ (use $Z_w=(x-w)-(x-z)$ and $Z_w^2=a+b-2(x-w)\cdot(x-z)$), and $-K_c+g_1+g_2+g_3+g_4=-K_h$. Then
\begin{align}
A_{\rm I.III}&=i\,(K_b-2K_a)\;f^{abd}\mathcal O_{\rm III}+{\rm c.c.}\qquad\text{(your Evolution\_three\_rho, with }-2K_a=2\KK(x'-w')\cdot\KK(y-z)\Phi(x)),\notag\\
B_{\rm I.III}&=i\,(-2K_b-2K_h+2K_a)\;f^{abd}\mathcal O_{\rm III}+{\rm c.c.},\notag\\
P_{\rm I.III}&=A+B=i\,(-K_b-2K_h)\;f^{abd}\mathcal O_{\rm III}+{\rm c.c.}
\end{align}
\textbf{Direct check of $P$.} Only two of your $\pp$ integrals reach $\vee$. The one with $\frac{[k^+(x-w)^2+\pp(x-z)^2]}{[\pp(x-z)^2-k^+(x-w)^2]}\frac{1}{\pp k^+}$ tends to $\frac{1}{\pp k^+}$ and gives $-K_b$. The one with $\frac{[k^+(x-w)^m-\pp(x-z)^m]}{[k^+(x-w)^2-\pp(x-z)^2]}\frac{1}{k^+-\pp}[\delta_{ik}\delta_{jm}-\dots]$ tends to $-\frac{1}{\pp}\frac{(x-z)^m}{(x-z)^2}\delta_{ik}\delta_{jm}$ and gives $-2K_h$. This is the same $P$, so your $\pp$ integration is correct (also checked by numerical integration).

\textbf{Fix P1.} The region-$P$ comparison (Sec.~\ref{sec:checks}) requires $P_{\rm I.III}=i(-K_b+2K_h)$. The term $-K_b$ is right, but the $\bar{\mathbb B}_3^\dagger$ term over $\kp+\Lambda<\pp<\vee$ has the wrong sign at large $\pp$. In the comparison, the structures with $K_h$ match the independent computation exactly once the sign is reversed, and they are off by exactly a factor $-1$ otherwise.

\textbf{Fix F3 carried over.} In the previous note the $\Phi$ terms of $A$ ($-2iK_a$) were replaced by $\mathcal D$,
\[
\mathcal D=-2i\int\KK(x'-w')\cdot\KK(y-w)\,K(x,w,z)\,f^{abd}\big[U^{ab'}(x')-U^{ab'}(w')\big]\rho^{b'}(x')\big[U^{de}(x)U^{bc}(w)-U^{de}(x)U^{bc}(y)\big]\rho^e(x)\rho^c(y)+{\rm c.c.}
\]
(the soft gluon exchanged between the measured gluon and the rotated charge $x$). Hence, corrected,
\begin{gather}
A_{\rm I.III}=iK_b\,f^{abd}\mathcal O_{\rm III}+\mathcal D,\qquad P_{\rm I.III}=i(-K_b+2K_h)\,f^{abd}\mathcal O_{\rm III}\quad(+{\rm c.c.}),\notag\\
\boxed{B_{\rm I.III}=i(-2K_b+2K_h)\,f^{abd}\mathcal O_{\rm III}-\mathcal D\quad(+{\rm c.c.})}
\end{gather}
The terms $\pm K_a$ (the edge kernel $\Phi$) are absent. In your notes they come from the poles at $\pp=\kp$ and multiply $\log\frac{k^+}{\Lambda}=\log\frac{\vee}{\Lambda}-\log\frac{\vee}{k^+}$.

\subsubsection{Row IV}
Prefactor $\frac{ig^4}{8\pi^5k^+}\to2i$. Colour operator
\[
\mathcal O_{\rm IV}=\big[U^{ab'}(x')\rho^{b'}(x')-U^{ab'}(w')\rho^{b'}(x')\big]\big[f^{cdc'}U^{ac}(w)U^{bd}(z)U^{be}(y)\rho^e(y)\rho^{c'}(x)-f^{abc}U^{be}(y)U^{cd}(x)\rho^e(y)\rho^d(x)\big].
\]
Your result has $L_P(a,b)\,\frac{(x'-w')^i(y-z)^j}{(x'-w')^2(y-z)^2}J_1^{ij}$ ($B=1$), $-\frac12\log\frac{\vee}{\Lambda+k^+}\frac{(x'-w')\cdot(y-z)}{(x'-w')^2(y-z)^2(z-w)^2}$ ($B=1$) and the blue $[\log\frac{\vee}{\Lambda}+\log(1-\frac{k^+}{\vee})]K_{\mathfrak B}$ ($A=1$). Therefore
\begin{gather}
A_{\rm I.IV}=2i\,K_{\mathfrak B}\,\mathcal O_{\rm IV}+{\rm c.c.},\qquad P_{\rm I.IV}=2i\,K^{P}_{\mathfrak B}\,\mathcal O_{\rm IV}+{\rm c.c.},\notag\\
\boxed{B_{\rm I.IV}=2i\,K^{L}_{\mathfrak B}\,\mathcal O_{\rm IV}+{\rm c.c.}}
\end{gather}
$P$ agrees with the large-$\pp$ limit of your integrand, $\frac{1}{k^+}\big[\frac{(x-z)^j(z-w)^i}{(z-w)^2}-\frac{(x-z)^j(x-w)^i}{2(x-w)^2}\big]\frac{1}{\pp(x-z)^2}$. Row IV is correct.

\subsubsection{Row V}
Here $0<\pp<\kp$, so $P_{\rm I.V}=0$. All your logarithms are $\log\frac{\vee}{\Lambda}+\log\frac{k^+-\Lambda}{\vee}=\log\frac{k^+-\Lambda}{\Lambda}$ ($A=1$, $B=-1$). Hence, for parts 1, 2a and 2b,
\begin{equation}
\boxed{B_{\rm I.V}=-A_{\rm I.V}}
\end{equation}
with $A_{\rm I.V}$ the $\log(\vee/\Lambda)$ terms of your Evolution files:
\begin{align}
B_{\rm I.V,1}={}&+\frac i2\,\NP\int\Big[e^{-ik\cdot(w'-z)}\KK(x'-w')\cdot\KK(y-z)\,\KK(z-x)\cdot\KK(y'-x)\notag\\
&\qquad\qquad+e^{-ik\cdot(w'-x)}\KK(x'-w')\cdot\KK(y'-x)\,\KK(z-x)\cdot\KK(y-z)\Big]\notag\\
&\times\big[U^{ab'}(x')-U^{ab'}(w')\big]\rho^{b'}(x')\Big[f^{adc}U^{db}(x)U^{ce}(z)\{\rho^e(y),\rho^b(y')\}-f^{abc}U^{cd}(y)U^{be}(y')\{\rho^d(y),\rho^e(y')\}\Big]+{\rm c.c.},
\end{align}
and $B_{\rm I.V,2}$ is minus Eq.~(4.6) of the previous note (parts 2a and 2b together, with $\NN\to\NP$).

%------------------------------------------------------------------------------------------
\subsection{Group II (no c.c.)}
\subsubsection{Row I}
\paragraph{Part 1 (three-$\rho$).} The two terms have the structure of Group I Row IV, with prefactors $-\frac{ig^4}{16\pi^5k^+}\to-i$ and $+i$:
\begin{align}
B_{\rm II.I,1}={}&-i\,\NP\int K^L_{\mathfrak B}\,\mathcal O^{(1)}+i\,\NP\int\bar K^L_{\mathfrak B}\,\mathcal O^{(2)},\notag\\
\mathcal O^{(1)}={}&\Big[U^{ac'}(w')\{\rho^{c'}(x'),\rho^b(y)\}-U^{bc'}(z)U^{ad'}(x')U^{c'e'}(y)\{\rho^{d'}(x'),\rho^{e'}(y)\}\Big]\notag\\
&\times\Big[f^{cbd}U^{ac}(w)\rho^d(x)-f^{acd}U^{bc}(z)U^{de}(x)\rho^e(x)\Big],\notag\\
\mathcal O^{(2)}={}&\Big[f^{c'bd'}U^{ac'}(w')\rho^{d'}(x)-f^{ac'd'}U^{bc'}(z)U^{d'e'}(x)\rho^{e'}(x)\Big]\notag\\
&\times\Big[U^{ac}(w)\{\rho^c(x'),\rho^b(y)\}-U^{bc}(z)U^{ad}(x')U^{ce}(y)\{\rho^d(x'),\rho^e(y)\}\Big].
\end{align}
$A$ and $P$ are the same with $K^L\to K$ and $K^L\to K^P$.

\paragraph{Part 2 (two-$\rho$).} Colour operator (your reduced form)
\begin{align*}
\mathcal O_{\rm II.I}={}&f^{c'bd'}f^{cbd}U^{ac'}(w')U^{ac}(w)\rho^{d'}(x')\rho^d(x)-f^{c'bd'}f^{acd}U^{ac'}(w')U^{bc}(z)U^{de}(x)\rho^{d'}(x')\rho^e(x)\\
&-f^{ac'd'}f^{cbd}U^{ac}(w)U^{bc'}(z)U^{e'd'}(x')\rho^{e'}(x')\rho^d(x)+N_cU^{e'd}(x')U^{de}(x)\rho^{e'}(x')\rho^e(x).
\end{align*}
With $a,b,a',b'$ as in Table~\ref{tab:L} and $D=a'b-ab'$, the $\log(\vee/k^+)$ content of your six parts 2a--2f is:
\begin{itemize}
\item 2a: $L_P(a,b)-L_P(a',b')$, so $B=1-1=0$;
\item 2b: blue $\log\frac{\vee-k^+}{\Lambda}$ ($B=0$; it gives $A$), and $-\frac{b}{aD}L_P(a,b)+\frac{b'}{a'D}L_P(a',b')$, so $B=\frac{-a'b+ab'}{aa'D}=-\frac{1}{aa'}$ times $K_B$;
\item 2c: $\frac{a'}{b'D}L_P(a',b')-\frac{a}{bD}L_P(a,b)$, so $B=\frac{a'b-ab'}{bb'D}=\frac{1}{bb'}$ times $K_C$;
\item 2d, 2e, 2f: the coefficients of $\log\frac{\vee}{\Lambda+k^+}$ and of the two $L_P$ cancel, $B=0$.
\end{itemize}
Here $K_B$ and $K_C$ are your brackets multiplying $\frac{(\pp+\kp)^2}{(\pp)^2}$ and $\frac{(\pp+\kp)^2}{(\kp)^2}$. They factorize:
\begin{align}
\frac{K_B}{aa'}&=\KK(x-w)\cdot\KK(x'-w')\,\Big[\KK(w-z)-\tfrac12\KK(x-z)\Big]\cdot\Big[\KK(w'-z)-\tfrac12\KK(x'-z)\Big],\notag\\
\frac{K_C}{bb'}&=K(x,x',z)\,\Big[\KK(z-w)-\tfrac12\KK(x-w)\Big]\cdot\Big[\KK(z-w')-\tfrac12\KK(x'-w')\Big].
\end{align}
With the prefactor $\frac{g^4}{8\pi^5k^+}\to2$:
\begin{equation}
A_{\rm II.I,2}=2\int\frac{K_B}{aa'}\,\mathcal O_{\rm II.I},\qquad P_{\rm II.I,2}=2\int\frac{K_C}{bb'}\,\mathcal O_{\rm II.I},\qquad
\boxed{B_{\rm II.I,2}=2\int\Big[\frac{K_C}{bb'}-\frac{K_B}{aa'}\Big]\mathcal O_{\rm II.I}}
\end{equation}
In $A$ the softer gluon $z$ is emitted by the measured gluon (or by $x$, with weight $-\frac12$). In $P$ the measured gluon is emitted by the harder gluon $z$ (or by $x$, with weight $-\frac12$). The two are related by exchanging the roles of $z$ and $w$, $w'$.

\subsubsection{Rows II and III}
Only four-$\rho$ terms $\propto\log(\vee/\Lambda)$: \textbf{$B=0$}.

\subsubsection{Row IV}\label{sec:II.IV}
Colour operator $\mathcal O_{\rm II.IV}(p',p)=\big[U^{e'c}(p')-U^{e'c}(x')\big]\big[U^{ec}(p)-U^{ce}(x)\big]\rho^{e'}(x')\rho^e(x)$. Your blue factor is $\log\frac{\vee}{\Lambda}+\log\frac{k^+}{\vee}$ ($A=1$, $B=-1$). The $\xi$ integral $\int_1^{k^+/\vee}\frac{d\xi}{\xi^2}$ hides a $+\log(\vee/k^+)$ at its lower end $\xi\to k^+/\vee$, i.e.\ $\pp\to\vee$. It is easiest to read off before the $w$, $w'$ and $z$ integrations. For $\pp\gg\kp$ the product of your two brackets is $(\frac{\pp+\kp}{\kp})^2\delta_{mm'}\delta_{kk'}$, times $\frac{\pp\kp}{(\pp+\kp)^4}$. This gives $\frac{1}{\kp\pp}$ with $y,y'\to z$. Hence ($\frac{g^4N_c}{8\pi^4k^+}\to2\pi N_c$ and $\frac{g^4N_c}{8\pi^5k^+}\to2N_c$)
\begin{align}
A_{\rm II.IV}&=2N_c\int_{x,x'}\Big[\int_zK(w,w',z)\Big]\KK(x-w)\cdot\KK(x'-w')\,\mathcal O_{\rm II.IV}(w',w),\notag\\
P_{\rm II.IV}&=2N_c\int_{x,x',z}K(w,w',z)\,K(x,x',z)\,\mathcal O_{\rm II.IV}(z,z),\notag\\
B_{\rm II.IV}&=P_{\rm II.IV}-A_{\rm II.IV}.
\end{align}
Here $\int_zK(w,w',z)=\pi[\frac1{\bar\epsilon}-\log(4\pi^2\mu^2(w-w')^2)]$ is your bracket. In $P$ both measured gluons are emitted by the harder gluon at $z$, which comes from $x$ and $x'$.

%------------------------------------------------------------------------------------------
\subsection{Group III (with c.c.)}
The overall sign of Rows III, V and VI is that of \texttt{Evolution\_three\_rho.tex} ($-\frac{4}{(2\pi)^3}$, $-\frac{2}{(2\pi)^3}$), which is opposite to \texttt{Group\_III.tex}. This applies to the $\log(\vee/k^+)$ terms as well, and to part 2 of Row III (fix F5).

\subsubsection{Rows I and II}
Same structure as Group I Row IV, with prefactor $-\frac{ig^4}{8\pi^5k^+}\to-2i$:
\begin{equation}
B_{\rm III.I(II)}=-2i\,\NP\int K^L_{\mathfrak B}\,\mathcal O_{\rm III.I(II)}+{\rm c.c.},
\end{equation}
with
\begin{align*}
\mathcal O_{\rm III.I}&=\big[U^{ab'}(x')U^{c'd'}(y)U^{c'b}(z)\rho^{b'}(x')\rho^{d'}(y)-U^{bd'}(z)U^{d'e'}(y)U^{ac'}(w')\rho^{c'}(x')\rho^{e'}(y)\big]\\
&\quad\times\big[f^{cbd}U^{ac}(w)\rho^d(x)-f^{acd}U^{bc}(z)U^{de}(x)\rho^e(x)\big],\\
\mathcal O_{\rm III.II}&=\big[U^{c'd'}(y)U^{ab'}(x')U^{bc'}(z)\rho^{d'}(y)\rho^{b'}(x')-U^{ab'}(x')\rho^b(y)\rho^{b'}(x')\big]\\
&\quad\times\big[f^{cbd}U^{ac}(w)\rho^d(x)-f^{acd}U^{bc}(z)U^{de}(x)\rho^e(x)\big].
\end{align*}

\subsubsection{Row III}\label{sec:G3III}
\paragraph{Part 1 (three-$\rho$).} Your result is the plus-prescription integral $\mathcal I_1$ plus the blue term $[\log\frac{\vee}{\Lambda}+\log(1-\frac{k^+}{\vee})]$. The blue term has $A=1$, $B=0$, and gives
\[
A_{\rm III.III,1}=-i\,\NN\int e^{-ik\cdot(y'-w)}\,K(y',x,z)\,\KK(x'-y')\cdot\KK(y-w)\;f^{c'd'a}\mathcal O_4\mathcal O_5+{\rm c.c.}
\]
with
\begin{align*}
\mathcal O_4&=U^{e'c'}(y')U^{d'b}(z)\rho^{e'}(x')-U^{bd'}(z)U^{c'e'}(x')\rho^{e'}(x'),\\
\mathcal O_5&=U^{ac}(w)\{\rho^c(y),\rho^b(x)\}-U^{bc}(z)U^{ad}(y)U^{ce}(x)\{\rho^d(y),\rho^e(x)\}.
\end{align*} The integral $\mathcal I_1=\int_0^{\vee-k^+}\frac{d\pp}{2\pi}(\dots)$ contains two hidden $\log(\vee/k^+)$:
\begin{enumerate}
\item the $[\frac1\pp]_+$ term with its phase gives $-\log(\vee/k^+)\times A_{\rm III.III,1}$ (Eq.~\eqref{eq:phase});
\item the term $-\frac1{k^+}\delta_{im}\delta_{jk'}$ behaves as $\frac{1}{k^+\pp}$ in the variable $w'$, where $y'=z+\xi(w'-z)$, $\xi=\frac{k^+}{k^++\pp}$, $d^2y'=\xi^2d^2w'$, and the phase becomes $e^{-ik\cdot(w'-w)}$. This gives $+\log(\vee/k^+)\times P_{\rm III.III,1}$ with
\[
P_{\rm III.III,1}=-i\,\NP\int e^{-ik\cdot(w'-w)}\,K(x',x,z)\,\KK(w'-z)\cdot\KK(y-w)\;f^{c'd'a}\,\mathcal O_4\big|_{y'\to z}\,\mathcal O_5+{\rm c.c.}
\]
\end{enumerate}
So $B_{\rm III.III,1}=P_{\rm III.III,1}-A_{\rm III.III,1}$. In $P$ the measured gluon $w'$ is emitted by the harder gluon $z$, which is exchanged between $x'$ and $x$.

\paragraph{Part 2 (two-$\rho$).} In the same way, from the blue term and the plus-prescription integral (with the sign of fix F5; prefactor $\frac{g^4}{8\pi^5k^+}\to2$):
\begin{align}
A_{\rm III.III,2}&=-2\,\NN\int e^{-ik\cdot(y'-w)}\,\KK(x-w)\cdot\KK(x'-y')\Big[\KK(z-w)\cdot\KK(y'-z)+\tfrac12K(x,y',z)\Big]\mathcal O_6\,\mathcal O_7+{\rm c.c.},\notag\\
P_{\rm III.III,2}&=-2\,\NP\int e^{-ik\cdot(w'-w)}\,K(x',x,z)\,\KK(w'-z)\cdot\Big[\KK(z-w)-\tfrac12\KK(x-w)\Big]\mathcal O_6\big|_{y'\to z}\,\mathcal O_7+{\rm c.c.},
\end{align}
with $\mathcal O_6=f^{c'd'a}[U^{e'c'}(y')U^{d'b}(z)-U^{bd'}(z)U^{c'e'}(x')]\rho^{e'}(x')$ and $\mathcal O_7=f^{cbd}U^{ac}(w)\rho^d(x)-f^{acd}U^{bc}(z)U^{de}(x)\rho^e(x)$. $P$ comes from the term $-\frac{\pp}{k^{+2}}[\frac{(x-z)^{k'}(z-w)^m}{(z-w)^2}-\frac{(x-z)^{k'}(x-w)^m}{2(x-w)^2}]$ of your bracket, which in the variable $w'$ behaves as $\frac{1}{k^+\pp(x-z)^2}[\dots]$. $B_{\rm III.III,2}=P-A$.

\subsubsection{Row IV}
Only four-$\rho$ terms $\propto\log(\vee/\Lambda)$: \textbf{$B=0$}.

\subsubsection{Rows V and VI}
Both have the longitudinal structure of Row III part 1, with twice its prefactor:
\begin{align}
A_{\rm III.V(VI)}&=-2i\,\NN\int e^{-ik\cdot(y'-w)}\,K(y',x,z)\,\KK(x'-y')\cdot\KK(y-w)\;\mathcal C_{\rm V(VI)}+{\rm c.c.},\notag\\
P_{\rm III.V(VI)}&=-2i\,\NP\int e^{-ik\cdot(w'-w)}\,K(x',x,z)\,\KK(w'-z)\cdot\KK(y-w)\;\mathcal C_{\rm V(VI)}\big|_{y'\to z}+{\rm c.c.},
\end{align}
with $\mathcal C_{\rm V}=f^{c'ca}[U^{e'c'}(y')\rho^{e'}(x')-U^{c'e'}(x')\rho^{e'}(x')][U^{ce}(x)U^{ad}(y)\rho^e(x)\rho^d(y)-U^{ad}(w)U^{ce}(x)\rho^e(x)\rho^d(y)]$ and $\mathcal C_{\rm VI}=f^{c'd'a}[U^{e'c'}(y')U^{d'b}(z)\rho^{e'}(x')-U^{bd'}(z)U^{c'e'}(x')\rho^{e'}(x')][U^{bc}(z)U^{ad}(y)U^{ce}(x)\rho^d(y)\rho^e(x)-U^{ad}(y)\rho^d(y)\rho^b(x)]$. $B=P-A$.

%------------------------------------------------------------------------------------------
\subsection{Summary of the rows}
\begin{center}\small
\begin{tabular}{@{}l l l l@{}}
row & $\log(\vee/k^+)$ terms $B$ & as written & status\\ \hline
I.I, II.II, II.III, III.IV, Eq.~(4rho) & 0 & -- & correct\\
I.II part 1 & $-A$ & explicit & F6 (bracket)\\
I.II part 2 & $-A$ & explicit & correct (F4 affects $A$ only)\\
I.III & $P-A$ & explicit & P1 ($+2K_h$), F3 ($\mathcal D$)\\
I.IV & $2iK^L_{\mathfrak B}$ & explicit & correct\\
I.V & $-A$ & explicit & correct\\
II.I parts 1, 2 & $P-A$ & explicit & correct\\
II.IV & $P-A$ & $-A$ explicit, $P$ hidden & correct\\
III.I, III.II & $-2iK^L_{\mathfrak B}$ & explicit & correct\\
III.III part 1, III.V, III.VI & $P-A$ & hidden & correct (sign of Evolution\_three\_rho)\\
III.III part 2 & $P-A$ & hidden & F5 (sign)\\
\end{tabular}
\end{center}
